\newpage
\section*{Part B: Exercise 7.1}

%---------------------------------------------------------------
\subsection*{Exercise 7.1, Question 1}
\textbf{Question.} Find the distance between the following pairs of points:\\
(i) $(2,3),\ (4,1)$\\
(ii) $(-5,7),\ (-1,3)$\\
(iii) $(a,b),\ (-a,-b)$

\textbf{Solution.} By the distance formula, the distance between $P(x_1,y_1)$ and $Q(x_2,y_2)$ is
\[PQ=\sqrt{(x_2-x_1)^2+(y_2-y_1)^2}.\]
\distnote

\textbf{(i)} $(x_1,y_1)=(2,3)$ and $(x_2,y_2)=(4,1)$:
\begin{align*}
\dist{d}{2}{3}{4}{1}{2\sqrt2}
\end{align*}
Here $\sqrt8=\sqrt{4\times2}=2\sqrt2$.

\textbf{(ii)} $(x_1,y_1)=(-5,7)$ and $(x_2,y_2)=(-1,3)$:
\begin{align*}
\dist{d}{-5}{7}{-1}{3}{4\sqrt2}
\end{align*}
Here $\sqrt{32}=\sqrt{16\times2}=4\sqrt2$.

\textbf{(iii)} $(x_1,y_1)=(a,b)$ and $(x_2,y_2)=(-a,-b)$:
\begin{align*}
d&=\sqrt{(-a-a)^2+(-b-b)^2}\\
&=\sqrt{(-2a)^2+(-2b)^2}\\
&=\sqrt{4a^2+4b^2}=\sqrt{4(a^2+b^2)}=2\sqrt{a^2+b^2}
\end{align*}
\slip{$\sqrt{4a^2+4b^2}$ is not $2a+2b$: the square root of a sum is not the sum of the square roots. Factor out $4$ first.}

\textbf{Verification.} For (iii) take $a=3$, $b=4$. The points are $(3,4)$ and $(-3,-4)$:
\[
\sqrt{(-3-3)^2+(-4-4)^2}=\sqrt{(-6)^2+(-8)^2}=\sqrt{36+64}=\sqrt{100}=10,
\]
and the formula gives $2\sqrt{3^2+4^2}=2\sqrt{25}=2\times5=10$. The two agree.

\ans{(i) $2\sqrt2$ units; (ii) $4\sqrt2$ units; (iii) $2\sqrt{a^2+b^2}$ units.}

%---------------------------------------------------------------
\subsection*{Exercise 7.1, Question 2}
\textbf{Question.} Find the distance between the points $(0,0)$ and $(36,15)$. Can you now find the distance between the two towns $A$ and $B$ discussed in Section 7.2.

(Section 7.2: A town $B$ is located 36 km east and 15 km north of the town $A$.)

\textbf{Solution.} By the distance formula, $d=\sqrt{(x_2-x_1)^2+(y_2-y_1)^2}$, with $(x_1,y_1)=(0,0)$ and $(x_2,y_2)=(36,15)$:
\begin{align*}
\dist{d}{0}{0}{36}{15}{39}
\end{align*}
Here $36^2=1296$, $15^2=225$, $1296+225=1521$ and $39^2=1521$, so $\sqrt{1521}=39$.

Now the towns. Take town $A$ as the origin $(0,0)$, one unit as 1~km, east as the positive $x$-direction and north as the positive $y$-direction. Then town $B$, 36 km east and 15 km north of $A$, is at $(36,15)$. So the distance between the towns is the distance between $(0,0)$ and $(36,15)$.

\ans{The distance between $(0,0)$ and $(36,15)$ is $39$ units. Yes, the distance between the towns $A$ and $B$ is $39$ km.}

%---------------------------------------------------------------
\subsection*{Exercise 7.1, Question 3}
\textbf{Question.} Determine if the points $(1,5)$, $(2,3)$ and $(-2,-11)$ are collinear.

\textbf{Solution.} Let $A(1,5)$, $B(2,3)$, $C(-2,-11)$.

Three points are collinear if the largest of the three distances equals the sum of the other two. By the distance formula $d=\sqrt{(x_2-x_1)^2+(y_2-y_1)^2}$:

\distnote
\begin{align*}
\distn{AB}{1}{5}{2}{3}\approx2.24\\
\dist{BC}{2}{3}{-2}{-11}{2\sqrt{53}\approx14.56}\\
\distn{AC}{1}{5}{-2}{-11}\approx16.28
\end{align*}
Here $\sqrt{212}=\sqrt{4\times53}=2\sqrt{53}$.

The largest distance is $AC\approx16.28$. The sum of the other two is
\[AB+BC\approx2.24+14.56=16.80\neq16.28\approx AC.\]
For an exact check, suppose $\sqrt5+\sqrt{212}=\sqrt{265}$. Squaring both sides, with $(p+q)^2=p^2+2pq+q^2$:
\[5+212+2\sqrt{5\times212}=265\ \Rightarrow\ 2\sqrt{1060}=48\ \Rightarrow\ \sqrt{1060}=24\ \Rightarrow\ 1060=24^2=576,\]
which is false. So $AB+BC\neq AC$.
\slip{Compare the largest distance with the sum of the \emph{other two}. Comparing $AB+AC$ with $BC$ proves nothing.}

\textbf{Alternative solution (section formula).} If $C$ lay on line $AB$, it would divide $AB$ in some ratio $k:1$, and by the section formula $\left(\dfrac{m_1x_2+m_2x_1}{m_1+m_2},\dfrac{m_1y_2+m_2y_1}{m_1+m_2}\right)$ with $A(1,5)$, $B(2,3)$, $m_1:m_2=k:1$,
\[C=\left(\frac{k(2)+1(1)}{k+1},\frac{k(3)+1(5)}{k+1}\right)=\left(\frac{2k+1}{k+1},\frac{3k+5}{k+1}\right).\]
From the $x$-coordinate: $\dfrac{2k+1}{k+1}=-2\Rightarrow 2k+1=-2k-2\Rightarrow k=-\dfrac34$.\\
From the $y$-coordinate: $\dfrac{3k+5}{k+1}=-11\Rightarrow 3k+5=-11k-11\Rightarrow k=-\dfrac87$.\\
The two values of $k$ differ, so no single ratio puts $C$ on line $AB$.

\ans{The points are not collinear.}

%---------------------------------------------------------------
\subsection*{Exercise 7.1, Question 4}
\textbf{Question.} Check whether $(5,-2)$, $(6,4)$ and $(7,-2)$ are the vertices of an isosceles triangle.

\textbf{Solution.} Let $A(5,-2)$, $B(6,4)$, $C(7,-2)$.

By the distance formula $d=\sqrt{(x_2-x_1)^2+(y_2-y_1)^2}$:

\distnote
\begin{align*}
\distn{AB}{5}{-2}{6}{4}\\
\distn{BC}{6}{4}{7}{-2}\\
\dist{AC}{5}{-2}{7}{-2}{2}
\end{align*}
(Here $\sqrt{37}\approx6.08$, and $\sqrt4=2$.)

\textbf{Do the points form a triangle?} They do if the sum of any two distances is greater than the third:
\begin{align*}
AB+BC&\approx6.08+6.08=12.16>AC=2\\
AB+AC&\approx6.08+2=8.08>BC\approx6.08\\
BC+AC&\approx6.08+2=8.08>AB\approx6.08
\end{align*}
\textbf{Is it isosceles?} A triangle is isosceles if two of its sides are equal. Here $AB=BC=\sqrt{37}$.

\ans{Yes, the points are the vertices of an isosceles triangle, with $AB=BC=\sqrt{37}$ units.}

%---------------------------------------------------------------
\subsection*{Exercise 7.1, Question 5}
\textbf{Question.} In a classroom, 4 friends are seated at the points $A$, $B$, $C$ and $D$ as shown in Fig.\ 7.8. Champa and Chameli walk into the class and after observing for a few minutes Champa asks Chameli, ``Don't you think $ABCD$ is a square?'' Chameli disagrees. Using distance formula, find which of them is correct.

\begin{center}
\ncertgrid{\fill (3,4) circle (3pt) node[above left]{$A$};
\fill (6,7) circle (3pt) node[above left]{$B$};
\fill (9,4) circle (3pt) node[above right]{$C$};
\fill (6,1) circle (3pt) node[above left]{$D$};}\\
{\small Fig.\ 7.8 (redrawn)}
\end{center}

\textbf{Solution.} \emph{Reading the figure.} In Fig.\ 7.8 the horizontal axis is labelled Columns and the vertical axis is labelled Rows. We take the column number as the $x$-coordinate and the row number as the $y$-coordinate. From the figure:
$A$ is in column 3, row 4, so $A(3,4)$; $B(6,7)$; $C(9,4)$; $D(6,1)$.

A quadrilateral is a square if all four sides are equal and both diagonals are equal. By the distance formula $d=\sqrt{(x_2-x_1)^2+(y_2-y_1)^2}$:

\distnote
\begin{align*}
\dist{AB}{3}{4}{6}{7}{3\sqrt2}\\
\dist{BC}{6}{7}{9}{4}{3\sqrt2}\\
\dist{CD}{9}{4}{6}{1}{3\sqrt2}\\
\dist{DA}{6}{1}{3}{4}{3\sqrt2}\\
\dist{AC}{3}{4}{9}{4}{6}\\
\dist{BD}{6}{7}{6}{1}{6}
\end{align*}
Here $\sqrt{18}=\sqrt{9\times2}=3\sqrt2$.

All four sides are equal ($3\sqrt2$ each) and the diagonals $AC$ and $BD$ are equal ($6$ each).
\slip{Four equal sides alone also describe a rhombus. Always check the diagonals before calling it a square.}

\textbf{Alternative solution.} By the converse of the Pythagoras theorem, $AB^2+BC^2=18+18=36=AC^2$ (the numbers $18$, $18$ and $36$ are those under the square roots above), so $\angle B=90^\circ$. A quadrilateral with four equal sides and one right angle is a square.

\ans{$ABCD$ is a square. Champa is correct and Chameli is wrong.}

%---------------------------------------------------------------
\subsection*{Exercise 7.1, Question 6}
\textbf{Question.} Name the type of quadrilateral formed, if any, by the following points, and give reasons for your answer:\\
(i) $(-1,-2),\ (1,0),\ (-1,2),\ (-3,0)$\\
(ii) $(-3,5),\ (3,1),\ (0,3),\ (-1,-4)$\\
(iii) $(4,5),\ (7,6),\ (4,3),\ (1,2)$

\textbf{Solution.}

\textbf{Method for every part.}
\begin{enumerate}[label=\arabic*.,leftmargin=*,itemsep=1pt,topsep=2pt]
\item Name the points $A,B,C,D$ in the order given.
\item Use the distance formula $d=\sqrt{(x_2-x_1)^2+(y_2-y_1)^2}$ to find the four sides $AB,BC,CD,DA$ and the two diagonals $AC,BD$.
\item Compare the lengths and read the type from the table.
\end{enumerate}

\begin{center}
\begin{tabular}{|p{0.6\linewidth}|l|}
\hline
\textbf{What the lengths show} & \textbf{Type of quadrilateral}\\ \hline
All four sides equal and diagonals equal & Square\\
All four sides equal but diagonals not equal & Rhombus\\
Opposite sides equal and diagonals equal (adjacent sides not equal) & Rectangle\\
Opposite sides equal but diagonals not equal & Parallelogram\\
Three of the four points lie on one straight line & No quadrilateral is formed\\ \hline
\end{tabular}
\end{center}
\slip{Check the order of the points. The quadrilateral is joined $A\to B\to C\to D\to A$, so $AC$ and $BD$ are the diagonals.}

\distnote

\textbf{(i)} $A(-1,-2)$, $B(1,0)$, $C(-1,2)$, $D(-3,0)$.
\begin{align*}
\dist{AB}{-1}{-2}{1}{0}{2\sqrt2}\\
\dist{BC}{1}{0}{-1}{2}{2\sqrt2}\\
\dist{CD}{-1}{2}{-3}{0}{2\sqrt2}\\
\dist{DA}{-3}{0}{-1}{-2}{2\sqrt2}\\
\dist{AC}{-1}{-2}{-1}{2}{4}\\
\dist{BD}{1}{0}{-3}{0}{4}
\end{align*}
(Here $\sqrt8=2\sqrt2$.) All four sides are equal ($2\sqrt2$) and the diagonals are equal ($4$). So $ABCD$ is a square.

\textbf{(ii)} $A(-3,5)$, $B(3,1)$, $C(0,3)$, $D(-1,-4)$.
\begin{align*}
\dist{AB}{-3}{5}{3}{1}{2\sqrt{13}}\\
\distn{BC}{3}{1}{0}{3}\\
\distn{AC}{-3}{5}{0}{3}
\end{align*}
(Here $\sqrt{52}=\sqrt{4\times13}=2\sqrt{13}$.) Then $AC+CB=\sqrt{13}+\sqrt{13}=2\sqrt{13}=AB$, so $A$, $C$, $B$ are collinear.

Check with the mid-point formula $\left(\dfrac{x_1+x_2}{2},\dfrac{y_1+y_2}{2}\right)$: the mid-point of $AB$ is $\left(\dfrac{-3+3}{2},\dfrac{5+1}{2}\right)=(0,3)=C$. So $C$ lies on $AB$.

Among any four points, if three are collinear, they cannot be three corners of a quadrilateral, whatever the order of the vertices. So no quadrilateral is formed.

\textbf{(iii)} $A(4,5)$, $B(7,6)$, $C(4,3)$, $D(1,2)$.
\begin{align*}
\distn{AB}{4}{5}{7}{6}\\
\dist{BC}{7}{6}{4}{3}{3\sqrt2}\\
\distn{CD}{4}{3}{1}{2}\\
\dist{DA}{1}{2}{4}{5}{3\sqrt2}\\
\dist{AC}{4}{5}{4}{3}{2}\\
\dist{BD}{7}{6}{1}{2}{2\sqrt{13}}
\end{align*}
(Here $\sqrt{18}=3\sqrt2$ and $\sqrt{52}=2\sqrt{13}$.) Opposite sides are equal: $AB=CD=\sqrt{10}$ and $BC=DA=3\sqrt2$. Adjacent sides are unequal ($\sqrt{10}\neq3\sqrt2$), so it is not a rhombus or a square. The diagonals are unequal ($AC=2$, $BD=2\sqrt{13}$), so it is not a rectangle. So $ABCD$ is a parallelogram.

\textbf{Verification for (iii).} By the mid-point formula, the mid-point of $AC$ is $\left(\dfrac{4+4}{2},\dfrac{5+3}{2}\right)=(4,4)$ and the mid-point of $BD$ is $\left(\dfrac{7+1}{2},\dfrac{6+2}{2}\right)=(4,4)$. The diagonals bisect each other, which confirms a parallelogram.

\ans{(i) $ABCD$ is a square. (ii) No quadrilateral is formed ($A$, $C$, $B$ are collinear). (iii) $ABCD$ is a parallelogram.}

%---------------------------------------------------------------
\subsection*{Exercise 7.1, Question 7}
\textbf{Question.} Find the point on the $x$-axis which is equidistant from $(2,-5)$ and $(-2,9)$.

\textbf{Solution.} A point on the $x$-axis has the form $(x,0)$ (its $y$-coordinate is $0$). Let $P(x,0)$, $A(2,-5)$, $B(-2,9)$. Since $PA=PB$, we have $PA^2=PB^2$.

By the distance formula $d^2=(x_2-x_1)^2+(y_2-y_1)^2$:
\[
PA^2=(2-x)^2+(-5-0)^2,\qquad PB^2=(-2-x)^2+(9-0)^2 .
\]
Since $(2-x)^2=(x-2)^2$ and $(-2-x)^2=(x+2)^2$, equate them and expand with $(a\pm b)^2=a^2\pm2ab+b^2$:
\begin{align*}
(x-2)^2+(-5)^2&=(x+2)^2+9^2\\
(x^2-4x+4)+25&=(x^2+4x+4)+81\\
x^2-4x+29&=x^2+4x+85\\
-4x+29&=4x+85\\
-8x&=56\\
x&=-7
\end{align*}
So $P=(-7,0)$.
\slip{The point is on the $x$-axis, so its $y$-coordinate is $0$ (not its $x$-coordinate). Take care with the sign in $(x+2)^2$.}

\textbf{Verification.} With $P(-7,0)$:
\begin{align*}
\distn{PA}{-7}{0}{2}{-5}\\
\distn{PB}{-7}{0}{-2}{9}
\end{align*}
Both distances are $\sqrt{106}$.

\ans{The required point is $(-7,0)$.}

%---------------------------------------------------------------
\subsection*{Exercise 7.1, Question 8}
\textbf{Question.} Find the values of $y$ for which the distance between the points $P(2,-3)$ and $Q(10,y)$ is 10 units.

\textbf{Solution.} By the distance formula $PQ=\sqrt{(x_2-x_1)^2+(y_2-y_1)^2}$ with $(x_1,y_1)=(2,-3)$, $(x_2,y_2)=(10,y)$ and $PQ=10$:
\begin{align*}
\sqrt{(10-2)^2+(y-(-3))^2}&=10\\
\sqrt{8^2+(y+3)^2}&=10\\
64+(y+3)^2&=100\qquad(\text{squaring both sides})\\
(y+3)^2&=36\\
y+3&=\pm6
\end{align*}
So $y+3=6$ or $y+3=-6$, that is $y=3$ or $y=-9$.
\slip{$(y+3)^2=36$ has two solutions, $+6$ and $-6$. Writing only $y+3=6$ loses the second point.}

\textbf{Verification.} For $y=3$, $Q=(10,3)$:
\[PQ=\sqrt{(10-2)^2+(3-(-3))^2}=\sqrt{8^2+6^2}=\sqrt{64+36}=\sqrt{100}=10.\]
For $y=-9$, $Q=(10,-9)$:
\[PQ=\sqrt{(10-2)^2+(-9-(-3))^2}=\sqrt{8^2+(-6)^2}=\sqrt{64+36}=\sqrt{100}=10.\]

\ans{$y=3$ or $y=-9$.}

%---------------------------------------------------------------
\subsection*{Exercise 7.1, Question 9}
\textbf{Question.} If $Q(0,1)$ is equidistant from $P(5,-3)$ and $R(x,6)$, find the values of $x$. Also find the distances $QR$ and $PR$.

\textbf{Solution.} Since $Q$ is equidistant from $P$ and $R$, $QP=QR$, so $QP^2=QR^2$.

By the distance formula $d^2=(x_2-x_1)^2+(y_2-y_1)^2$:
\begin{align*}
QP^2&=(5-0)^2+(-3-1)^2=5^2+(-4)^2=25+16=41\\
QR^2&=(x-0)^2+(6-1)^2=x^2+5^2=x^2+25
\end{align*}
Equating, $x^2+25=41$, so $x^2=16$ and $x=4$ or $x=-4$.

\textbf{The distances.} $QR^2=x^2+25=16+25=41$ for both values of $x$, so $QR=\sqrt{41}$.

For $PR$, take $(x_1,y_1)=(5,-3)$ and $(x_2,y_2)=(x,6)$:
For $x=4$, $R=(4,6)$:
\begin{align*}
\distn{PR}{5}{-3}{4}{6}
\end{align*}
For $x=-4$, $R=(-4,6)$:
\begin{align*}
\dist{PR}{5}{-3}{-4}{6}{9\sqrt2}
\end{align*}
(Here $\sqrt{162}=\sqrt{81\times2}=9\sqrt2$.)
\slip{$x^2=16$ gives both $x=4$ and $x=-4$, and the two points $R$ give different values of $PR$.}

\ans{$x=4$ or $x=-4$; $QR=\sqrt{41}$ units; $PR=\sqrt{82}$ units when $x=4$ and $PR=9\sqrt2$ units when $x=-4$.}

%---------------------------------------------------------------
\subsection*{Exercise 7.1, Question 10}
\textbf{Question.} Find a relation between $x$ and $y$ such that the point $(x,y)$ is equidistant from the point $(3,6)$ and $(-3,4)$.

\textbf{Solution.} Let $P(x,y)$ be equidistant from $A(3,6)$ and $B(-3,4)$, so $PA=PB$ and $PA^2=PB^2$.

By the distance formula $d^2=(x_2-x_1)^2+(y_2-y_1)^2$:
\[
PA^2=(x-3)^2+(y-6)^2,\qquad PB^2=(x-(-3))^2+(y-4)^2=(x+3)^2+(y-4)^2 .
\]
Equate them and expand with $(a\pm b)^2=a^2\pm2ab+b^2$:
\begin{align*}
(x-3)^2+(y-6)^2&=(x+3)^2+(y-4)^2\\
(x^2-6x+9)+(y^2-12y+36)&=(x^2+6x+9)+(y^2-8y+16)\\
x^2+y^2-6x-12y+45&=x^2+y^2+6x-8y+25\\
-6x-12y+45&=6x-8y+25\\
45-25&=6x+6x+12y-8y\\
20&=12x+4y\\
3x+y&=5\qquad(\text{dividing by }4)
\end{align*}

\textbf{Verification.} The point $(1,2)$ satisfies $3x+y=3(1)+2=5$. Its distances from $A(3,6)$ and $B(-3,4)$ are
\begin{align*}
\dist{PA}{3}{6}{1}{2}{2\sqrt5}\\
\dist{PB}{-3}{4}{1}{2}{2\sqrt5}
\end{align*}
Both equal $\sqrt{20}$, so the point is equidistant.

\ans{The required relation is $3x+y=5$.}
